\documentclass[11pt,a4paper]{article}

\usepackage[utf8]{inputenc}
\usepackage[T1]{fontenc}
\usepackage{lmodern}
\usepackage{microtype}
\usepackage[margin=2.5cm]{geometry}
\usepackage{graphicx}
\usepackage{amsmath,amssymb}
\usepackage{listings}
\usepackage{xcolor}
\usepackage{hyperref}
\usepackage{booktabs}
\usepackage{caption}
\usepackage{subcaption}
\usepackage{enumitem}

\hypersetup{
  colorlinks=true,
  linkcolor=black,
  urlcolor=blue,
  citecolor=black
}

% MiniZinc listing style
\definecolor{mznKw}{RGB}{0,90,160}
\definecolor{mznCom}{RGB}{120,120,120}
\definecolor{mznStr}{RGB}{170,55,55}

\lstdefinelanguage{minizinc}{
  morekeywords={
    include, par, var, array, of, int, bool, set, constraint,
    forall, exists, sum, alldifferent, solve, satisfy, minimize, maximize,
    if, then, else, endif, in, div, mod, where, output, predicate, function
  },
  sensitive=true,
  morecomment=[l]{\%},
  morestring=[b]"
}

\lstset{
  language=minizinc,
  basicstyle=\ttfamily\footnotesize,
  keywordstyle=\color{mznKw}\bfseries,
  commentstyle=\color{mznCom}\itshape,
  stringstyle=\color{mznStr},
  numbers=left,
  numberstyle=\tiny\color{gray},
  numbersep=8pt,
  breaklines=true,
  showstringspaces=false,
  frame=single,
  framesep=4pt,
  rulecolor=\color{gray!40},
  captionpos=b,
  xleftmargin=10pt
}

\title{Robot Movement Planning\\\large Constraint Programming Project Report}
\author{Alessio Russo (mat. 376856)}
\date{}

\begin{document}
\maketitle

% =============================================================================
\section{Problem Statement}
\label{sec:problem}

We consider a rectangular layer of boxes arranged on a discrete grid. Each box
has dimensions \(x \times y\), where the \(x\)-axis is aligned with the belt
direction and the \(y\)-axis is orthogonal to it. The target packaging layer
contains \(n\) boxes, whose destination positions are denoted by
\((P_x[i], P_y[i])\), for \(i = 1,\dots,n\). These positions are assumed to be
non-overlapping and to tile a rectangular grid.

A robotic gripper with two parallel plates, each of size \(S\), can pick a
contiguous sequence of boxes along the \(x\)-axis. The total physical length of
the sequence must be at most \(S + x\), which corresponds to a maximum number
of boxes equal to
\[
    \left\lfloor \frac{S+x}{x} \right\rfloor .
\]
This formulation accounts for a maximum admissible overflow of \(x/2\) on each
side of the gripper.

\begin{figure}[h]
    \centering
    \includegraphics[width=0.8\textwidth]{images/pick_and_release.pdf}
    \caption{Example of a picking and releasing operations.}
    \label{fig:pick_and_release_example}
\end{figure}

At each release operation, only one gripper plate can open. Therefore, the
opening side must face a free direction, either the external boundary of the
layer or a grid position that has not yet been filled. This constraint couples the
geometric construction of groups with the temporal order in which they are
released.

Figure~\ref{fig:pick_and_release_example} shows an example of the
picking and releasing operations. The problem can therefore be decomposed
into two related subproblems:

\begin{enumerate}[label=(\Alph*)]
  \item \textbf{Grouping.} Partition the \(n\) boxes into groups that are
        contiguous along the \(x\)-axis, belong to the same row, and fit within
        the gripper capacity.
  \item \textbf{Scheduling.} Determine an ordering of the groups and assign an
        opening plate side to each group, so that the opening plate always
        faces a free direction at the time of release.
\end{enumerate}

\subsection{Parameters and Derived Constants}

Based on this formulation, the MiniZinc model uses the following input
parameters:

\begin{itemize}
  \item the total number of boxes in the layer, denoted by
        \texttt{n\_boxes};
  \item the grid dimensions of the layer, denoted by
        \texttt{boxes\_along\_x} and \texttt{boxes\_along\_y};
  \item the box dimension parallel to the gripper, denoted by
        \texttt{boxes\_x\_dim}. In the natural orientation,
        \texttt{boxes\_x\_dim} corresponds to \(x\); if the layer is rotated,
        it corresponds to \(y\);
  \item the gripper size \(S\), denoted by \texttt{grip\_size}.
\end{itemize}

These parameters are declared as follows:
\vspace{0.5cm}
\begin{lstlisting}
par int: n_boxes;
par int: boxes_x_dim;
par int: boxes_along_x;
par int: boxes_along_y;
par int: grip_size;
\end{lstlisting}
\vspace{0.5cm}
From these inputs, the model derives three constants. First, the maximum number
of boxes that can be picked in a single operation is computed directly from the
gripper geometry:
\[
    \texttt{max\_pkble\_boxes}
    =
    \left\lfloor
    \frac{\texttt{grip\_size} + \texttt{boxes\_x\_dim}}
         {\texttt{boxes\_x\_dim}}
    \right\rfloor .
\]
Second, the minimum number of groups required to cover one row is obtained by
ceiling division:
\[
    \texttt{groups\_per\_row}
    =
    \left\lceil
    \frac{\texttt{boxes\_along\_x}}
         {\texttt{max\_pkble\_boxes}}
    \right\rceil .
\]
Finally, the total number of groups is obtained by multiplying the number of
groups per row by the number of rows:
\[
    \texttt{n\_groups}
    =
    \texttt{boxes\_along\_y}
    \cdot
    \texttt{groups\_per\_row}.
\]

The corresponding declarations are:
\vspace{0.5cm}
\begin{lstlisting}
par int: max_pkble_boxes = (grip_size + boxes_x_dim) div boxes_x_dim;

par int: groups_per_row = (boxes_along_x + max_pkble_boxes - 1)
                           div max_pkble_boxes;

par int: n_groups = boxes_along_y * groups_per_row;
\end{lstlisting}
% =============================================================================
\section{Grouping}
\label{sec:model}
This part of the model admits several possible encodings. A natural but
expensive formulation is a \emph{fully enumerated} encoding, where each group
explicitly stores the identifiers of all boxes it contains. Since groups may
have different lengths, this representation usually requires a two-dimensional
array with padding values for shorter groups.
\vspace{0.5cm}
\begin{lstlisting}
array[1..n_groups, 1..max_pkble_packs] of var 0..n_packs: groups;
\end{lstlisting}
\vspace{0.5cm}

Here, the value \(0\) is used as a dummy value to represent an empty slot.
This representation makes group membership explicit and easy to inspect.
However, it introduces a large number of decision variables and requires
additional constraints to enforce the structure of a valid solution. In
particular, the model must ensure that each box appears exactly once, that
non-empty positions inside each group are ordered, and that the groups are
consistent with the geometric ordering of the layer.

%\vspace{0.5cm}
%\begin{lstlisting}
%% Each pack identifier 1..n_packs appears exactly once.
%constraint
%  forall(v in 1..n_packs) (
%    count(array1d(groups), v) = 1
%  );
%
%% Non-empty entries are placed before padding values.
%constraint
%  forall(g in 1..n_groups, j in 1..max_pkble_packs-1) (
%    groups[g,j] = 0 -> groups[g,j+1] = 0
%  );
%
%% Pack identifiers are strictly increasing inside each group.
%constraint
%  forall(g in 1..n_groups, j in 1..max_pkble_packs-1) (
%    groups[g,j] != 0 /\ groups[g,j+1] != 0 ->
%      groups[g,j] < groups[g,j+1]
%  );
%\end{lstlisting}
%\vspace{0.5cm}

A second possibility is a \emph{binary assignment} encoding, where a Boolean
variable states whether pack \(i\) belongs to group \(g\). This representation
can be written as follows:

\vspace{0.5cm}
\begin{lstlisting}
array[1..n_groups, 1..n_packs] of var bool: belongs_to_group;
\end{lstlisting}
\vspace{0.5cm}

In this case, group membership is represented by the truth value of
\texttt{belongs\_to\_group[g,i]}. Although this encoding is expressive, it is
even larger than the fully enumerated representation. It also does not encode
contiguity or same-row membership directly. These properties must be
reconstructed through additional channeling constraints.

For this reason, the adopted model represents each group using its
\emph{start index} and \emph{length}. An equivalent formulation could use the
start and end indices instead. This interval-style representation is more
compact because each group is described by only two main variables. Contiguity
is implicit: once the start index and the length are known, the boxes in the
group are uniquely determined. The same-row constraint can also be expressed
directly by requiring the interval not to cross a row boundary. As a result,
the model avoids explicit membership variables and reduces the number of
constraints needed to describe valid groups.

\vspace{0.5cm}
\begin{lstlisting}
array[1..n_groups] of var 1..n_packs: group_start;
array[1..n_groups] of var 1..max_pkble_packs: group_len;
\end{lstlisting}
\vspace{0.5cm}

The adopted encoding uses two arrays: \texttt{group\_start}, which stores the
first box of each group, and \texttt{group\_len}, which stores the number of
boxes in that group.

\vspace{0.5cm}
\begin{lstlisting}
constraint group_start[1] = 1;

constraint forall(g in 1..n_groups - 1)(
  group_start[g + 1] = group_start[g] + group_len[g]
);

constraint group_start[n_groups] + group_len[n_groups] - 1 = n_packs;

constraint forall(g in 1..n_groups)(
  (group_start[g] - 1) div packs_along_x =
  (group_start[g] + group_len[g] - 2) div packs_along_x
);
\end{lstlisting}
\vspace{0.5cm}

The first three constraints define the groups as consecutive intervals over
the row-major ordering of the packs. The first group starts from pack \(1\),
and each following group starts immediately after the last pack of the previous
group. The final constraint in this block forces the last group to end at pack
\(n\). Together, these conditions ensure that all packs are covered exactly
once, without gaps or overlaps.

The fourth constraint enforces the geometric feasibility of each group. Since
the gripper can only pick boxes aligned along the \(x\)-axis, a group cannot
cross from one row to the next. The constraint therefore compares the row index
of the first pack in the group with the row index of the last pack. If the two
indices are equal, the whole interval lies on the same row.

\section{Scheduling}

\begin{figure}[t!]
  \centering
  \includegraphics[width=0.65\linewidth]{images/sample_pallet.pdf}
  \caption{Pallet layout produced by the grouping stage: 51 groups
  tiling the $11 \times 17$ grid, each colored
  distinctly. Groups alternate $3$--$4$--$4$ across each row.}
  \label{fig:sample-pallet}
\end{figure}

The second part of the model determines the release order of the groups and
selects which gripper plate opens for each release. In this case, the most
direct representation consists in assigning to each group a single release
position. The array \texttt{drop\_order} stores the time step at which each
group is released, while \texttt{plate} stores the plate opened for that release.
The value \(0\) denotes the left plate, and the value \(1\) denotes the right
plate.

\vspace{0.5cm}
\begin{lstlisting}
array[1..n_groups] of var 1..n_groups: drop_order;
array[1..n_groups] of var 0..1: plate;
\end{lstlisting}
\vspace{0.5cm}

The only global constraint required to make \texttt{drop\_order} a valid
schedule is \texttt{alldifferent}. Since every value in \texttt{drop\_order}
belongs to the domain \(1,\dots,\texttt{n\_groups}\), forcing all values to be
different ensures that each release step is assigned to exactly one group.

\vspace{0.5cm}
\begin{lstlisting}
constraint alldifferent(drop_order);
\end{lstlisting}
\vspace{0.5cm}

The feasibility of a release depends on the position of the adjacent groups in the same
vertical column. As shown in Figure~\ref{fig:sample-pallet}, groups are
numbered from left to right within each row, and then from top to
bottom across rows.
Therefore, since each row contains \texttt{groups\_per\_row} groups,
two groups that are aligned in consecutive rows have indices
that differ by \texttt{groups\_per\_row}. For a generic group \(g\),
the vertical neighborhood is therefore given by 
$g \pm \texttt{groups\_per\_row},$ whenever these indices are within the valid range.
Groups in the first pallet row have no upper neighbor,
while groups in the last pallet row have no lower neighbor.

The final constraint encodes the physical restriction on the gripper plates:
when a group is released, the selected plate must open toward a free direction.
A direction is free if the group lies on the corresponding pallet boundary, or
if the group in that direction is scheduled to be released later and is
therefore still absent from the pallet at the current release step.

Since the relevant neighbors are above and below the current group, the
constraint is written by separating the first row, the internal rows, and the
last row. This avoids invalid array accesses at the pallet boundaries.

\vspace{0.5cm}
\begin{lstlisting}
constraint forall(g in 1..n_groups where g <= groups_per_row)(
  plate[g] = 0 \/
  (
    plate[g] = 1 /\
    drop_order[g + groups_per_row] > drop_order[g]
  )
);

constraint forall(g in 1..n_groups where
  g > groups_per_row /\ g <= n_groups - groups_per_row
)(
  (
    plate[g] = 0 /\
    drop_order[g - groups_per_row] > drop_order[g]
  )
  \/
  (
    plate[g] = 1 /\
    drop_order[g + groups_per_row] > drop_order[g]
  )
);

constraint forall(g in 1..n_groups where g > n_groups - groups_per_row)(
  (
    plate[g] = 0 /\
    drop_order[g - groups_per_row] > drop_order[g]
  )
  \/
  plate[g] = 1
);
\end{lstlisting}
\vspace{0.5cm}

The first constraint handles groups in the first pallet row. These groups have
no upper neighbor, so opening plate \(0\) is always feasible. Opening plate
\(1\), instead, is feasible only if the lower neighbor is released later.

The second constraint handles the internal rows. In this case, both vertical
neighbors exist. Opening plate \(0\) is feasible only if the upper neighbor
\(g-\texttt{groups\_per\_row}\) is released later. Similarly, opening plate
\(1\) is feasible only if the lower neighbor
\(g+\texttt{groups\_per\_row}\) is released later.

The third constraint handles groups in the last pallet row. These groups have
no lower neighbor, so opening plate \(1\) is always feasible. Opening plate
\(0\) is feasible only if the upper neighbor is released later.

%An equivalent and more compact formulation is to precompute the vertical
%neighborhood using parameter arrays \texttt{upper\_neighbor} and
%\texttt{lower\_neighbor}. The sentinel value \(0\) denotes the absence of a
%neighbor at a pallet boundary.
%
%\vspace{0.5cm}
%\begin{lstlisting}
%array[1..n_groups] of par int: upper_neighbor =
%  [ if g > groups_per_row
%    then g - groups_per_row
%    else 0
%    endif
%  | g in 1..n_groups ];
%
%array[1..n_groups] of par int: lower_neighbor =
%  [ if g <= n_groups - groups_per_row
%    then g + groups_per_row
%    else 0
%    endif
%  | g in 1..n_groups ];
%\end{lstlisting}
%\vspace{0.5cm}
%
%Using these arrays, the release constraint can be written more compactly. A
%group can open plate \(0\) if it has no upper neighbor, or if its upper neighbor
%is released later. The same condition applies symmetrically to plate \(1\) and
%the lower neighbor.
%
%\vspace{0.5cm}
%\begin{lstlisting}
%constraint forall(g in 1..n_groups)(
%  (
%    plate[g] = 0 /\
%    if upper_neighbor[g] = 0 then
%      true
%    else
%      drop_order[upper_neighbor[g]] > drop_order[g]
%    endif
%  )
%  \/
%  (
%    plate[g] = 1 /\
%    if lower_neighbor[g] = 0 then
%      true
%    else
%      drop_order[lower_neighbor[g]] > drop_order[g]
%    endif
%  )
%);
%\end{lstlisting}
%\vspace{0.5cm}
%
%This formulation expresses the same release rule as the row-separated version,
%but separates the neighborhood computation from the scheduling condition. It
%also makes the constraint shorter and easier to inspect. The explicit
%row-separated formulation is preferable when avoiding conditional expressions,
%whereas the neighbor-array formulation provides a more compact representation
%of the same logic.

The model is a pure feasibility problem:

\vspace{0.5cm}
\begin{lstlisting}
solve satisfy;
\end{lstlisting}
\vspace{0.5cm}

No objective function is required. Any assignment of \texttt{drop\_order} and
\texttt{plate} that satisfies the vertical-neighbor constraint defines a valid
pick-and-place plan.

\begin{figure}[h]
  \centering
  \begin{subfigure}[t]{0.24\linewidth}
    \includegraphics[width=\linewidth]{images/scheduling/schedule_step_01.pdf}
  \end{subfigure}\hfill
  \begin{subfigure}[t]{0.24\linewidth}
    \includegraphics[width=\linewidth]{images/scheduling/schedule_step_10.pdf}
  \end{subfigure}\hfill
  \begin{subfigure}[t]{0.24\linewidth}
    \includegraphics[width=\linewidth]{images/scheduling/schedule_step_25.pdf}
  \end{subfigure}\hfill
  \begin{subfigure}[t]{0.24\linewidth}
    \includegraphics[width=\linewidth]{images/scheduling/schedule_step_50.pdf}
  \end{subfigure}
  \caption{Example of a scheduling. The red
  border marks the group being placed; the blue arrow shows the
  opening plate side.}
  \label{fig:sample-schedule}
\end{figure}

% =============================================================================
\section{Driving the model from Python}
\label{sec:python}

The Python layer is intentionally thin --- it is a wrapper around
\texttt{minizinc-python} that handles the geometry computations the model
deliberately leaves out, and produces visual artifacts of solutions. Three
scripts compose the pipeline:

\begin{description}
  \item[\texttt{sample.py}] Solves a single instance and renders the
        layout and a step-by-step drop schedule as PDFs. The Python side
        is responsible for picking the orientation that fits more boxes
        on the pallet (\texttt{lib/utils.py}: \texttt{gen\_pallet\_info}),
        deriving \texttt{boxes\_x\_dim} and the gripper size $S = 3x$,
        and translating the model output (the four arrays
        \texttt{group\_start}, \texttt{group\_len}, \texttt{drop\_order},
        \texttt{plate}) into a high-level \emph{schedule} of drop
        actions.
  \item[\texttt{benchmark.py}] Runs the same model across the cartesian
        product of five pallet presets and six pack sizes, against every
        available solver, with a 5-minute timeout. Each run is
        encapsulated and timing/error information is recorded in a
        single \texttt{benchmark/results.json} file.
  \item[\texttt{plot\_benchmark.py}] Consumes the JSON benchmark and
        produces the comparison plots referenced in
        Section~\ref{sec:benchmark}.
\end{description}

The renderer in \texttt{lib/visualize.py} is orientation-agnostic: it
consumes already-rotated cell dimensions (\texttt{cell\_w}, \texttt{cell\_h})
and the four solution arrays, and outputs both an overall pallet figure
(colored by group) and one figure per drop step (with an arrow showing the
opening plate side).

%\subsection{Worked example}
%\label{sec:example}
%
%To make the output of the model concrete, we run \texttt{sample.py} on a
%$70 \times 70$ pack and an $800 \times 1200$ pallet. The geometry helper
%chooses the rotation that maximises pack count, yielding an
%$11 \times 17$ grid (\(n = 187\) boxes). With \(S = 3 \cdot 70 = 210\),
%the gripper capacity is \(\lfloor (210 + 70)/70 \rfloor = 4\) boxes, and
%each row of 11 boxes is split into 3 groups of sizes $3+4+4$, for a
%total of \(n\_groups = 51\).
%
%Figure~\ref{fig:sample-schedule} samples four steps from the 51-step
%schedule produced by the model. The opening plate is highlighted by a
%red border on the active group and a blue arrow pointing in the
%direction the plate opens; the side it points to is always empty at the
%moment of release.
%
%
%The schedule starts in the bottom-right corner and works back toward
%the top-left: at step 1 the gripper drops \(G_{51}\) opening to the
%left (its left neighbor \(G_{50}\) has not yet been placed), and the
%final step places \(G_1\) opening to the left into the pallet edge
%(no left neighbor exists, so the boundary sentinel applies).

% =============================================================================
\section{Benchmarking}
\label{sec:benchmark}

\subsection{Setup}

The benchmark runs the MiniZinc model on every $(\text{pallet},
\text{pack})$ combination from the lists in \texttt{benchmark.py} (see
Table~\ref{tab:bench-grid}), against every solver registered with the
local MiniZinc installation, with a per-run timeout of 300 seconds.

\begin{table}[h]
  \centering
  \begin{tabular}{lll}
    \toprule
    \textbf{Pallet preset} & \textbf{Dimensions (mm)} \\
    \midrule
    Europallet         & $1200 \times 800$  \\
    Industrie          & $1200 \times 1000$ \\
    US 40$\times$48    & $1219 \times 1016$ \\
    Half-pallet        & $800  \times 600$  \\
    Australian         & $1165 \times 1165$ \\
    \bottomrule
  \end{tabular}
  \quad
  \begin{tabular}{l}
    \toprule
    \textbf{Pack size (mm)} \\
    \midrule
    $50 \times 50$   \\
    $70 \times 70$   \\
    $100 \times 100$ \\
    $120 \times 80$  \\
    $150 \times 150$ \\
    $200 \times 200$ \\
    \bottomrule
  \end{tabular}
  \caption{Pallet presets and pack sizes used in the benchmark.}
  \label{tab:bench-grid}
\end{table}

The gripper size is fixed at $S = 3 \cdot \texttt{boxes\_x\_dim}$, the
maximum allowed by the project specification.

\subsection{Overall scaling}

Figure~\ref{fig:scaling-overall} shows solve time as a function of the
total number of boxes on the pallet, on a logarithmic axis, for every
solver across all instances. The geometric mean and min/max range per
solver is summarised in Figure~\ref{fig:solver-summary}.

\begin{figure}[h]
  \centering
  \includegraphics[width=0.85\linewidth]{images/scaling_overall_n_boxes.pdf}
  \caption{Solve time vs. total number of boxes --- all solvers, all
  instances. Log scale on the time axis.}
  \label{fig:scaling-overall}
\end{figure}

\begin{figure}[h]
  \centering
  \includegraphics[width=0.7\linewidth]{images/solver_summary.pdf}
  \caption{Per-solver geometric mean of solve time (bar) with min/max
  whiskers across the full benchmark set.}
  \label{fig:solver-summary}
\end{figure}

\subsection{Per-pallet breakdown}

Splitting the data by pallet preset (Figure~\ref{fig:scaling-per-pallet})
isolates the effect of pack count from pallet shape: for a fixed pallet,
the curves are functions of pack count alone, while differences between
pallets reflect the impact of \texttt{boxes\_along\_x} and
\texttt{boxes\_along\_y} on the number of groups.

\begin{figure}[h]
  \centering
  \begin{subfigure}[t]{0.48\linewidth}
    \includegraphics[width=\linewidth]{images/scaling_europallet_n_boxes.pdf}
    \caption{Europallet}
  \end{subfigure}\hfill
  \begin{subfigure}[t]{0.48\linewidth}
    \includegraphics[width=\linewidth]{images/scaling_industrie_n_boxes.pdf}
    \caption{Industrie}
  \end{subfigure}

  \medskip

  \begin{subfigure}[t]{0.48\linewidth}
    \includegraphics[width=\linewidth]{images/scaling_us_40x48_n_boxes.pdf}
    \caption{US 40$\times$48}
  \end{subfigure}\hfill
  \begin{subfigure}[t]{0.48\linewidth}
    \includegraphics[width=\linewidth]{images/scaling_halfpallet_n_boxes.pdf}
    \caption{Half-pallet}
  \end{subfigure}

  \medskip

  \begin{subfigure}[t]{0.48\linewidth}
    \includegraphics[width=\linewidth]{images/scaling_australian_n_boxes.pdf}
    \caption{Australian}
  \end{subfigure}
  \caption{Solve time vs. number of boxes, broken down by pallet preset.}
  \label{fig:scaling-per-pallet}
\end{figure}

\subsection{Wall time vs.\ solve time}

The gap between Python-measured wall time and solver-reported solve time
captures MiniZinc flattening and inter-process overhead.
Figure~\ref{fig:wall-vs-solve} plots the two against each other; points
above the diagonal indicate runs in which the harness cost dominated the
search itself.

\begin{figure}[h]
  \centering
  \includegraphics[width=0.7\linewidth]{images/wall_vs_solve.pdf}
  \caption{Wall time vs. solve time per run, log--log. The diagonal
  represents zero overhead; all points lie above it.}
  \label{fig:wall-vs-solve}
\end{figure}

% =============================================================================
\section{Conclusion}
\label{sec:conclusion}

The two-stage formulation --- grouping the pallet into pickable runs, then
scheduling the drops with a plate side per group --- maps cleanly onto
MiniZinc decision variables and yields a compact, fully feasible model. The
key modelling choice is to fix the number of groups statically and force
the partition into a prefix-sum shape; this lets us pre-compute the
neighbor structure as parameters rather than variables, keeping the
plate-collision constraint a simple disjunction over indexed lookups.

The benchmark confirms that the model scales well across realistic pallet
sizes within the 5-minute budget, and exposes the expected per-solver
trade-offs between flattening overhead and search performance.

\end{document}
